\documentclass[11pt]{article}
\usepackage{preamble}
\setlength{\parskip}{4pt}

\begin{document}

\daybanner{AP Statistics \textperiodcentered\ Unit 3 \textperiodcentered\ Topic 3.14 \textperiodcentered\ B: 2026-12-18 \textperiodcentered\ E: 2027-02-01 \textperiodcentered\ TEACHER KEY}{Two Simulations for Mystery Tokens}

\sectionbanner{CED TETHER}

{\small
\begin{itemize}[leftmargin=1.5em,itemsep=2pt,topsep=2pt]
  \item \textbf{VAR-1.J} Identify questions suggested by variation between observed and expected counts in categorical data. \textbf{[Skill 1.A]}
  \item \textbf{VAR-1.J.1} Variation between what we find and what we expect to find may be random or not.
\end{itemize}
}
\textbf{Essential question:} What must a simulation vary to decide whether observed categorical discrepancies are surprising?\par
\textbf{DOK-3 skill:} 4.B \textperiodcentered\ \textbf{FRQ pattern:} {\small\texttt{fixed-\allowbreak{}count-\allowbreak{}vs-\allowbreak{}null-\allowbreak{}generation}}\par
\textbf{DOK rationale:} Part (c) coordinates the null model, what each simulation preserves, the chi-square tail, and an evidence threshold to adjudicate two designs and explain the false assurance produced by the rejected design.\par

\frameworkphaseheader{Do Now (first take) $\rightarrow$ Explore (video + follow-along) $\rightarrow$ Exit (finish + turn in)}{1 $\rightarrow$ 3}{45}{%
\begin{itemize}[leftmargin=1.5em,itemsep=2pt,topsep=2pt]
  \item Display the board and ask for one sentence using the study description.
  \item Begin the 8.1 video follow-along at minute 5.
  \item Return to the board at minute 33. Students finish the shortened parts and justify their judgment.
  \item Collect at the door. Score part (c) only using the three required elements.
\end{itemize}
}{%
\begin{itemize}[leftmargin=1.5em,itemsep=2pt,topsep=2pt]
  \item Write a one-sentence first take before the video.
  \item Complete the video follow-along worksheet.
  \item Finish (a)--(c), revise the first take in the exit line, and turn in the sheet.
\end{itemize}
}{%
\begin{itemize}[leftmargin=1.5em,itemsep=2pt,topsep=2pt]
  \item Which specific number or study detail supports your choice?
  \item What claim would go beyond the evidence, and why?
\end{itemize}
}{Circulate and ask for evidence. Accept a concise response; do not require omitted calculations or a second full inference write-up.}

\begin{callout}{\IconWarn\ WATCH FOR}{calloutred}
\begin{itemize}[leftmargin=1.5em,itemsep=2pt,topsep=2pt]
  \item Choosing a speaker without a reason.
  \item Copying numbers without explaining what they support.
\end{itemize}
\end{callout}

\sectionbanner{SCORING PART (c) --- E / P / I}

\textbf{Expected elements}
\begin{itemize}[leftmargin=1.5em,itemsep=2pt,topsep=2pt]
  \item \textbf{varied-counts} (required) --- Justifies Amara using new counts under the claimed probabilities
  \item \textbf{fixed-counts} (required) --- Explains that Leo redraws the same counts and cannot model their variation
  \item \textbf{bounded-conclusion} (required) --- Uses 0.186 > 0.05 to call the result unsurprising without claiming proof
\end{itemize}

{\small\renewcommand{\arraystretch}{1.25}
\noindent\begin{tabular}{|p{0.5in}|p{5.9in}|}
\hline
\textbf{E} & Explains both simulation mechanisms and uses the tail proportion for a bounded conclusion. \\ \hline
\textbf{P} & Shows substantive valid reasoning for at least one required element, but does not meet all E requirements. Credit correct reasoning even when another claim is wrong. \\ \hline
\textbf{I} & Shows no substantive valid reasoning for any required element; a name, unsupported choice, or copied number alone does not count. \\ \hline
\end{tabular}}

\clearpage
\focusbanner{TODAY'S PROBLEM --- annotated}

Suppose a company claims 10{,}000 mystery tokens are $50\%$ red, $30\%$ teal, and $20\%$ gold. An SRS of 80 has the shown counts.\par\smallskip \textbf{Leo:} ``Mix and redraw these same 80 tokens without replacement; recompute $\chi^2$ for 1,000 trials.''\par \textbf{Amara:} ``For each trial, generate 80 independent colors with probabilities 0.50, 0.30, and 0.20, then recompute $\chi^2$.'' In Amara's 1,000 trials, 186 statistics are at least the observed value.\par
\begin{center}
\textbf{Token sample}\par\smallskip
\begin{tabular}{|l|c|c|}
\hline
\textbf{Color} & \textbf{Claim} & \textbf{Observed} \\ \hline
\textbf{Red} & $0.50$ & 32 \\ \hline
\textbf{Teal} & $0.30$ & 30 \\ \hline
\textbf{Gold} & $0.20$ & 18 \\ \hline
\end{tabular}
\end{center}
\begin{firsttakebox}{FIRST TAKE (5 min)}
Which plan can produce a different number of red tokens on the next trial? Explain in one sentence.\par
\answer{Ungraded. Everyday language is enough before the video; formal vocabulary belongs in the finish.}
\end{firsttakebox}

\textit{Rules callout ``EXPECTED COUNTS AND SIMULATION (keep on the page)'' is printed on the student sheet.}\par\medskip

\dokbadge{1}~\textbf{(a)} Find the three expected color counts from the company's percentages.\par
\answer{$80(0.50)=40$ red, $80(0.30)=24$ teal, and $80(0.20)=16$ gold.}

\dokbadge{2}~\textbf{(b)} The observed statistic is $\chi^2=3.35$. Check the red contribution $(32-40)^2/40$, then compute the simulation proportion $186/1{,}000$.\par
\answer{The red contribution is $(32-40)^2/40=1.60$. The simulated upper-tail proportion is $186/1{,}000=0.186$.}

\dokbadge{3}~\textbf{(c)} In 3--4 sentences, choose a simulation and justify what varies from trial to trial. Use the simulated proportion to decide whether the result is unusual at a 5 percent cutoff, and explain why this does not prove the company's claim.\par
\answer{Amara is justified because each trial generates new counts using the claimed color probabilities. Leo redraws all of the same 80 tokens, so the counts and $\chi^2=3.35$ never change; his output cannot show sample-to-sample variation. Since $0.186>0.05$, the observed result is not unusual by the stated cutoff. This leaves chance variation plausible but does not prove the company's color claim.}

\begin{summaryexitbox}{EXIT}
One thing the video changed about my first take: \hrulefill
\end{summaryexitbox}

\end{document}
